% !TEX root = ../main.tex
\chapter{Introduction}
\label{ch:introduction}
\label{sec:intro}

Consider a decision you have to be able to defend two years from now.

A clinical operations team asks a language model system whether a compound is
approved for a paediatric indication. The system answers. Someone acts on the
answer. Later---in an audit, an inspection, a dispute---the question arrives:
what was that answer based on? Not whether the model is generally accurate. Which
document, in which version, as of which date, under which jurisdiction, supported
this specific claim.

Most current validation practice cannot answer that question, and the reason is
not that the practice is careless. It is that the practice measures the wrong
thing. The dominant approach evaluates a model by sampling questions with known
answers and scoring the outputs against them---\emph{reference-based evaluation},
in the terms of the natural-language-generation literature. It produces a
number---accuracy, or its complement, often reported as a hallucination
rate---and that number is then used as evidence of fitness for deployment.

The central argument of this book can be stated in one sentence. \textbf{In
regulated settings, the relevant question is not merely whether an answer is
accurate. It is whether the organization can reconstruct why this particular
assertion was produced, identify the authoritative and versioned evidence on which
it rests, establish that the evidence applies to the resolved case context, verify
the asserted relation, and show that reliance was authorized under the institution's
governing rules.} Accuracy is one input to that question and settles none of it.

This book argues that reference-based evaluation cannot support the use it is put
to,
for two reasons that compound.

The first is a classification problem. Pooling every wrong answer into one
metric conflates failures with different causes and different remedies. A system
that faithfully reports an erroneous source has a data problem. A system that
fabricates a plausible-looking figure has a generation problem. A system that
correctly reports a regulation that was superseded last year, or that applies in
a different jurisdiction, has a context-resolution problem. These are three
separate defects, addressed by three separate teams, on three separate budgets.
One number does not tell you which one to spend.

The second reason is worse, because it closes off the obvious remedy. The
system applies no canonical form to its inputs. It computes a conditional distribution over token sequences given
token sequences; its inputs are strings; and two semantically identical
queries---a contraposed conditional, a quantity written in digits rather than
words, an active rather than passive phrasing---are two distinct points in input
space. Nothing in the architecture guarantees they map to the same output, and
measurement confirms they frequently do not. A database normalizes before it
compares. A compiler canonicalizes before it optimizes. This system does
neither. Nor does the problem stop at paraphrase: the same input, re-asked in the
same session, can produce a different assertion as context accumulates. So a
passing test does not extend to materially identical requests, and does not
securely cover even the string that was tested.

The consequence for validation is direct. If a test case passes, you have
learned that one string produced an acceptable output on one occasion. You have
not learned that the semantically identical string will, or that the same string
will next week, or that the same string will later in the same conversation. A
passing test does not generalize to the paraphrase, so the test suite does not
certify the behavior it appears to certify. Accuracy is a property of a sample,
and this sample does not support the inference that validation requires of it.

That is a negative result, and on its own it would be unhelpful. The
constructive claim is that a different property does not have this defect.
\emph{Grounding}---whether an assertion is supported by evidence applicable to
the decision at hand, through a mechanism actually sensitive to that
evidence---is a property of the relation between an assertion and an evidential
base, not of a sampled output. It is stable under paraphrase in the way accuracy
is not, because it is a fact about the pipeline rather than about a draw from it.
It is therefore the property a governance regime can audit.

That claim is made precise later, in the sharpest form the argument takes: no
accuracy figure identifies or bounds the grounding rate over the unit square. No accuracy figure constrains the rate at which a system asserts
beyond its evidence, in either direction. A system can be highly accurate and
almost entirely ungrounded, and the second fact is the one that predicts what
happens when the questions change.

\paragraph{What this book is not arguing.} Not that these systems are
unreliable in a way that precludes use. I build them, in exactly the regulated
settings this book is about, and they work. Not that accuracy is uninteresting;
it is interesting, and it is the wrong thing to certify against. And not that
better models will not help. They will, on some axes and not others, and which
axes those are follows from the architecture itself.

\paragraph{Contributions.} Six, in the order they appear.

\begin{enumerate}
  \item \textbf{A redefinition of hallucination, indexed to a decision context.}
        Existing definitions ask whether an assertion is supported by a supplied
        source. This one asks the prior question---whether that source governs the
        case---by requiring support to be applicable at a resolved context of
        jurisdiction, valid time, material facts, and intended use, and by
        requiring the assertion to have been produced \emph{because of} the evidence
        rather than merely to agree with it. The consequence is a failure class that
        faithfulness and attribution metrics score as clean by construction, and a
        definition evaluable at the point of assertion without a reference answer.
  \item A five-category taxonomy following from it---\textsc{Sound},
        \textsc{Lucky}, \textsc{Error}, \textsc{Fabrication}, and
        \textsc{Misapplied}---in which the fifth, assertion grounded in authentic
        but inapplicable evidence, is a distinct control failure with its own
        remedy and owner.
  \item A proof that aggregate accuracy neither identifies nor bounds the grounding
        rate, so a validation program reporting accuracy alone cannot distinguish
        radically different reliance profiles.
  \item An architectural account of why the failures are structural, grounded in the
        specific architecture and training regime these systems inherit, and an
        argument that the absence of a canonical form makes accuracy uncertifiable
        rather than merely incomplete.
  \item A set of nineteen invariants separating what assertion behavior does not
        supply from what a system must be built to supply, each tied to a control.
  \item A reference architecture for governed retrieval, a definition of
        institutional warrant, and an account of why the division of ownership
        across data, policy, and platform functions is itself a principal cause of
        failure.
\end{enumerate}

\paragraph{Scope.} This book is about one class of system: models that emit
natural-language assertions by autoregressive sampling over a token vocabulary,
fitted by next-token likelihood on large text corpora, and optionally post-trained
on human preferences. That is the decoder-only transformer language model
established at scale by the GPT-3 generation and everything built on that pattern
since. It is not about deep learning generally. Generative adversarial networks and
diffusion models are out of scope, as are encoder-only classifiers, rankers, and
extractors, and the conventional predictive models that make up most of an
enterprise analytics estate. Those components assert nothing, so grounding is not a
coherent question about them---which is not to say they are easy to validate,
since they share several of the problems described here. The boundary is stated
precisely, with the reason it falls where it does, alongside the account of the
architecture.

\paragraph{Structure.} Chapter~\ref{ch:architecture-training} identifies the
systems in scope and derives, from their architecture and training regime, why the
failures this book describes are structural. Chapter~\ref{ch:taxonomy} redefines
hallucination as ungrounded assertion, gives the five categories that follow,
proves that aggregate accuracy does not identify the grounding rate, and defines
the resolved context and the seven dimensions of grounding --- along with why
retrieval and citation do not supply them.
Chapter~\ref{ch:canonical-form} shows that neither form invariance nor
repeatability holds, and what that costs a validation program.
Chapter~\ref{ch:invariants} turns those failures into a list: what these systems do
not supply, what a system must be built to supply instead, and why grounding rather
than accuracy is the property that can be measured.

The second half is constructive. Chapter~\ref{ch:governed-rag} gives a reference
architecture, Chapter~\ref{ch:scale-and-scope} says what in it is actually tunable,
Chapter~\ref{ch:warrant-tiers} defines warrant with reliance tiers and the
organizational analysis that goes with it, and Chapter~\ref{ch:evaluation} gives an
evaluation protocol. Chapter~\ref{ch:framework} then builds the control model that
generalizes the architecture --- five surfaces, a governance component, an algebra
for composing components --- and the three chapters after it govern the surfaces one
at a time: the prompt before the model runs, the output before it is delivered, and
the system as a whole over time. A closing chapter on related work positions the
argument.

\paragraph{Notation.} Formal notation is used where it makes a claim checkable
and avoided where it would only make one look rigorous. Two operators, five
defined categories, one proposition with a two-line proof. Every definition is
stated over observable behavior, so that any claim made here can be tested
against a running system.

\section{Premise: What Generation Does Not Establish}
\label{ch:epistemic-premise}
\label{sec:premise}

One premise carries the rest of the book, and it is deliberately modest. A
system's producing a statement is a different thing from that statement being
supported by evidence, applicable to the case at hand, checked, and permitted for
use. Any architect will grant the distinction in every other context: a report that
renders is not a report that reconciles, and a value that populates a field is not
a value that passed validation. The premise says only that generated text is
subject to the same distinction, and that the generation event is not what closes
it.

The argument below needs no position on what a model understands or represents
internally. It needs the distinction stated formally, which is the business of this
section.

\subsection{Notation}
\label{sec:notation}

\begin{center}
\begin{tabular}{@{}ll@{}}
\toprule
$M$ & the language model: a parameterized generative function \\
$S$ & the deployed system: $M$ together with its decoding, retrieval, policy,
      and control layers \\
$x$ & an input---a prompt, a structured request, a conversation state \\
$p, q$ & propositions or claims \\
$c$ & the resolved context of a decision \\
$E$ & an evidentiary basis \\
$V$ & a verification procedure \\
$I$ & an institution or accountable organization \\
\bottomrule
\end{tabular}
\end{center}

The distinction between $M$ and $S$ is load-bearing rather than pedantic. Every
property this book identifies as necessary for reliance is a property of $S$, and
none is a property of $M$. A model has no version of an audit record, an
authorization, or a resolved context; a system can have all three. Where the
literature reports a result about models, I say $M$; where the object of
governance is at issue, I say $S$.

\subsection{The resolved context}
\label{sec:resolved-context}

A great deal of enterprise information is true only relative to a governing frame.
Call that frame the \emph{resolved context} and write
\begin{equation}
  c = \{\, j,\ t_v,\ f,\ u \,\},
  \label{eq:context}
\end{equation}
where

\begin{center}
\begin{tabular}{@{}l>{\raggedright\arraybackslash}p{10.6cm}@{}}
\toprule
$j$ & jurisdiction or governing regime: the legal, regulatory, or policy authority
      whose rules apply \\
$t_v$ & \emph{valid time}: the effective date the assertion is about---the as-of date
      against which the governing rules are to be assessed \\
$f$ & material case facts: patient facts, customer facts, transaction facts,
      product facts, or whatever else is decision-relevant \\
$u$ & intended use: the decision class, the reliance context, what the answer will
      be used to do \\
\bottomrule
\end{tabular}
\end{center}

What these four have in common is what decides membership: each is something that
must be settled about the \emph{question} before the question is well posed. Ask
whether a filing may use the accelerated pathway and you have asked nothing
answerable until the regime, the as-of date, the case facts and the intended use are
fixed. That is why $c$ contains these and not other true and relevant things.

A second time dimension is required, and it fails that test, which is why it does not
belong in $c$. Write $t_k$ for \emph{transaction time}: the state of the evidential
base the system actually had access to---which corpus release, ingested to what date.
A requester does not supply $t_k$ and could not; it is a fact about the system at the
moment it answered, not about the question it was asked. The distinction is the
ordinary bitemporal one and it is not decoration.

$t_v$ is a property of the \emph{question}: which date's rules govern this decision.
$t_k$ is a property of the \emph{answer's provenance}: what the system could have
known when it answered. Correctness is assessed against $t_v$; reconstruction requires $t_k$; and their
relation is checkable in \emph{both} directions, which is the whole reason two
timestamps are needed:
\begin{align}
  t_k \prec t_v &\;\Longrightarrow\; \text{the base may omit instruments in force at }
    t_v \quad \text{(omission)}, \label{eq:bitemporal-stale}\\
  t_k \succ t_v &\;\Longrightarrow\; \text{the base may contain instruments not in
    force at } t_v \quad \text{(anachronism)}. \label{eq:bitemporal-anachronism}
\end{align}

Omission is the familiar case: a rule amended before $t_v$ but ingested after $t_k$
leaves an assertion ungrounded in a way no inspection of the retrieved evidence will
reveal, because the missing instrument is missing.

Anachronism is the other half and in regulated work it is the more dangerous of the two,
because the answer is confidently wrong rather than incompletely right. A corpus current
to today, answering a question as-of two years ago, will retrieve provisions that had
not yet been made. Applying a rule to a transaction that preceded it is a familiar and
serious error in its own right, and the system will produce a clean citation to a real
instrument while committing it.

Neither failure is prevented by having a recent corpus, and both are prevented by the
same condition. $\Applicable(E, c)$ must be read to require \emph{interval containment}
on the temporal component of $c$: the instrument's own validity period includes $t_v$.
That excludes the not-yet-enacted and the already-superseded together, and it is a
filter on the evidence rather than a property of the corpus. Temporal validity is
therefore part of applicability rather than a dimension of its own, which is why the
conditions on grounded evidence, set out later, do not list it separately. Both directions are also detectable at the aggregate level
without knowing what any amendment said: compare the corpus ingest watermark against
the decision date, per source class, and flag either sign of the difference. Systems that record only one timestamp cannot
perform the comparison, and most record only one.

$\Resolved(c)$ holds when the components of \eqref{eq:context} are explicit rather
than assumed or inferred without validation. $t_k$, for the reason above, is not
among them: it is recorded rather than resolved, and belongs to what the audit record
must preserve. The distinction between explicit and inferred is the whole of
it: a system that guesses the jurisdiction correctly has not resolved the context,
because nothing in the record establishes that the guess was checked, and an
auditor cannot distinguish a correct guess from a lucky one.

The truth predicate is correspondingly indexed. $\True(p \mid c)$ is truth relative
to a resolved context, and the relation that fails is:
\begin{equation}
  \True(p \mid j_1, t_1, f_1, u_1)
  \nentails
  \True(p \mid j_2, t_2, f_2, u_2).
  \label{eq:no-transfer}
\end{equation}

\subsection{The non-entailments}

Write $\Pr_M(\gen(p) \mid x)$ for the probability that the model $M$, given input
$x$, generates a token sequence expressing $p$. This is a generation event and nothing
more --- a fact about $M$'s output distribution. It is deliberately not called an
assertion: $\Assert_S(p)$, defined in Chapter~\ref{ch:taxonomy}, is a property of the
deployed system $S$ and is behavioral rather than probabilistic, for reasons given
there. The two are kept in separate notation because conflating them is the error this
section exists to block. The premise is that a high value of $\Pr_M(\gen(p) \mid x)$
does not establish any of the properties on which reliance depends:

\begin{align}
  \Pr_M(\gen(p) \mid x) \text{ high}
    &\nentails \True(p \mid c),            \label{eq:ne-true}\\
  \Pr_M(\gen(p) \mid x) \text{ high}
    &\nentails \Supports(E, p, c),         \label{eq:ne-supported}\\
  \Pr_M(\gen(p) \mid x) \text{ high}
    &\nentails \Applicable(E, c),          \label{eq:ne-applicable}\\
  \Pr_M(\gen(p) \mid x) \text{ high}
    &\nentails \Verified(V, p, c),         \label{eq:ne-verified}\\
  \Pr_M(\gen(p) \mid x) \text{ high}
    &\nentails \Warranted_I(p \mid c).     \label{eq:ne-warranted}
\end{align}

Compactly, and in the order in which an institution needs them:
\begin{equation}
  \Pr_M(\gen(p) \mid x) \text{ high}
  \nentails \True(p \mid c)
  \nentails \Supports(E, p, c)
  \nentails \Warranted_I(p \mid c).
  \label{eq:ne-chain}
\end{equation}

These are non-entailments, and the distinction matters because the claim is
frequently misread as stronger than it is. They do not assert that the properties
fail to co-occur. A model can and routinely does produce propositions that are
true, well supported, applicable, verifiable, and fit for reliance. The claim is
that the generation event is not what establishes those properties, so the
generation event cannot be what an institution points to when asked to defend
them. Each arrow in \eqref{eq:ne-chain} marks a gap that something other than the
model has to close.

Nor is the claim that the properties are unrelated. High generation probability is
correlated with truth on the distribution the model was fitted to, sometimes
strongly. Correlation on a training distribution is not entailment, and what concerns this
book is what happens when the distribution of interest is not that one.

\subsection{What follows}

The premise is weak, and that is its value: it can be granted without conceding
anything contentious, and it is sufficient for everything in this book.

Two clarifications, because both are easy to misread from the non-entailments
alone. The claim is not that model output is unreliable. Output is often excellent,
and the argument concerns what can be established about an individual assertion
rather than what can be expected of a population of them. And the claim is not that
the properties in \eqref{eq:ne-chain} are rare. In a well-built system most
assertions have all of them. The point is that they are established elsewhere, by
components that can be pointed to.

That is the whole of the constructive program. If generation does not establish
support, applicability, verification, or authorization, then those properties are
supplied by something in $S$ that is not $M$---and identifying, building, and
measuring that something is what the rest of this book is about.
